\documentclass[11pt]{article}

% ---- Page geometry (arxiv single-column) ----
\usepackage[margin=1in]{geometry}

% ---- Standard packages ----
\usepackage{amsmath,amssymb,amsthm}
\usepackage{mathtools}
\usepackage{algorithm}
\usepackage{algpseudocode}
\usepackage{booktabs}
\usepackage{multirow}
\usepackage{graphicx}
\usepackage{hyperref}
\usepackage{cleveref}
\usepackage{xcolor}
\usepackage{enumitem}

\usepackage{caption}
\usepackage{subcaption}
\usepackage{tikz}
\usetikzlibrary{arrows.meta,positioning,calc,patterns,decorations.pathreplacing}
\usepackage{pgfplots}
\pgfplotsset{compat=1.17}
\usepackage{listings}

% ---- Theorem environments ----
\newtheorem{theorem}{Theorem}[section]
\newtheorem{lemma}[theorem]{Lemma}
\newtheorem{proposition}[theorem]{Proposition}
\newtheorem{corollary}[theorem]{Corollary}
\newtheorem{definition}[theorem]{Definition}
\newtheorem{remark}[theorem]{Remark}
\newtheorem{example}[theorem]{Example}

% ---- Macros ----
\newcommand{\R}{\mathbb{R}}
\newcommand{\N}{\mathbb{N}}
\newcommand{\E}{\mathbb{E}}
\newcommand{\Prob}{\mathbb{P}}
\newcommand{\Var}{\mathrm{Var}}
\newcommand{\calS}{\mathcal{S}}
\newcommand{\calA}{\mathcal{A}}
\newcommand{\calO}{\mathcal{O}}
\newcommand{\calZ}{\mathcal{Z}}
\newcommand{\calH}{\mathcal{H}}
\newcommand{\calG}{\mathcal{G}}
\newcommand{\calP}{\mathcal{P}}
\newcommand{\calF}{\mathcal{F}}
\newcommand{\calC}{\mathcal{C}}
\newcommand{\calL}{\mathcal{L}}
\newcommand{\calM}{\mathcal{M}}
\newcommand{\calD}{\mathcal{D}}
\newcommand{\calT}{\mathcal{T}}
\newcommand{\calW}{\mathcal{W}}
\newcommand{\hb}{\rightarrow_{\mathrm{hb}}}
\newcommand{\conc}{\parallel}
\newcommand{\abs}[1]{\lvert #1 \rvert}
\newcommand{\norm}[1]{\lVert #1 \rVert}
\newcommand{\zo}{\mathsf{Zono}}
\newcommand{\marace}{\textsc{Marace}}
\DeclareMathOperator{\poly}{poly}
\DeclareMathOperator{\diag}{diag}
\DeclareMathOperator{\Lip}{Lip}
\DeclareMathOperator{\rank}{rank}
\DeclareMathOperator{\ESS}{ESS}

\title{\textbf{MARACE: Interaction Race Detection via Happens-Before-Aware\\Abstract Interpretation for Multi-Agent RL Safety Verification}}
\date{}

\begin{document}
\maketitle

\begin{abstract}
Multi-agent reinforcement learning (MARL) systems deployed in safety-critical domains
exhibit \emph{interaction races}: hazardous joint states that arise only under
specific relative execution orderings of concurrent agent policies. We formalize
interaction races through a happens-before (HB) partial-order semantics and prove
a \emph{separation theorem} showing that race detection is PSPACE-hard in general
but decidable in polynomial time when policies are piecewise-linear (ReLU) networks.
Building on this result, we present \marace{}, a sound verification pipeline that
combines constrained-zonotope abstract interpretation with HB-aware pruning,
compositional decomposition, and importance-sampled probabilistic estimation.
On two multi-agent benchmarks, \marace{} achieves perfect recall (1.00) across
all configurations with total analysis times under 0.52 seconds for up to 8 agents.
HB-aware pruning reduces the adversarial search space by $201\times$ compared
to na\"{i}ve exploration. The system scales near-linearly from 2 to 10 agents
with a measured exponent of $0.93$, enabling practical safety analysis of
deployed MARL systems.
\end{abstract}

% ============================================================
\section{Introduction}
\label{sec:intro}
% ============================================================

Multi-agent reinforcement learning has achieved impressive results in domains
ranging from autonomous driving~\cite{shalev2016safe} to warehouse logistics~\cite{li2020warehouse}
and robotic coordination~\cite{busoniu2008comprehensive}. As these systems move
toward deployment in safety-critical settings, a fundamental verification challenge
emerges: even when individual agent policies satisfy local safety properties, their
\emph{concurrent} execution may produce hazardous joint states under specific
relative timing and ordering of actions.

We call this phenomenon an \emph{interaction race}, by analogy with data races in
concurrent programming~\cite{lamport1978time,flanagan2009fasttrack}. Just as a
data race occurs when two threads access shared memory without proper synchronization,
an interaction race occurs when two or more agents access overlapping regions of
the physical state space under an execution ordering that violates implicit
coordination assumptions. Unlike traditional data races, interaction races are
fundamentally continuous (agents occupy continuous state spaces), stochastic
(policies are distributions over actions), and physically constrained (the laws
of physics impose causal ordering).

\paragraph{Motivating example.}
Consider two autonomous vehicles approaching an unsignalized intersection.
If Vehicle~A enters the intersection before Vehicle~B has observed A's position
update, B may proceed under the assumption that the intersection is clear.
Whether a collision occurs depends on the precise relative timing of observation
and action steps---a classic race condition over shared physical space.

\paragraph{Challenges.}
Three fundamental obstacles make interaction race detection difficult:
\begin{enumerate}[leftmargin=*,nosep]
\item \textbf{Non-projectability.} Interaction races are inherently multi-agent
phenomena that cannot be detected by analyzing individual policies in isolation
(\Cref{thm:non-proj}).
\item \textbf{Statistical intractability.} Under reasonable assumptions on
policy smoothness, the probability of observing a race in random simulation
decreases exponentially with the state-space dimension (\Cref{thm:stat-intract}).
\item \textbf{Computational hardness.} Exact interaction race detection for
general nonlinear policies is PSPACE-hard (\Cref{thm:pspace-hard}).
\end{enumerate}

\paragraph{Our approach.}
We overcome these obstacles through a combination of three techniques.
First, we establish that for \emph{piecewise-linear} (PWL) policies---such
as ReLU neural networks---interaction race detection is in P (\Cref{thm:pwl-p}).
Second, we develop a \emph{constrained-zonotope abstract domain} augmented with
happens-before ordering constraints that soundly over-approximates the set of
reachable joint states while eliminating physically infeasible execution orderings.
Third, we employ \emph{compositional decomposition} that partitions agents into
small interaction groups, enabling verification to scale near-linearly with
the number of agents.

\paragraph{Contributions.}
\begin{enumerate}[leftmargin=*,nosep]
\item We formalize the notion of \emph{interaction race} for MARL systems,
prove a separation theorem establishing its complexity (\Cref{sec:theory}),
and show that the problem admits efficient solutions for PWL policies.

\item We develop a \emph{happens-before-aware abstract interpretation} framework
using constrained zonotopes with HB pruning, including convergence guarantees
for the analysis of ReLU networks (\Cref{sec:abstract}).

\item We present \marace{}, a complete verification pipeline integrating
compositional decomposition, adversarial schedule search via MCTS, and
importance-sampled probability estimation (\Cref{sec:system}).

\item We evaluate \marace{} on highway intersection and warehouse corridor
benchmarks, demonstrating perfect recall, near-linear scaling to 10 agents,
and a $201\times$ search space reduction from HB pruning (\Cref{sec:eval}).
\end{enumerate}

% ============================================================
\section{Interaction Race Theory}
\label{sec:theory}
% ============================================================

We now develop the formal foundations for interaction races in multi-agent
systems. Throughout, we consider a system of $n$ agents with indices
$\calA = \{1,\ldots,n\}$ operating in a shared continuous state space
$\calS \subseteq \R^d$.

\subsection{Multi-Agent Execution Model}

\begin{definition}[Multi-agent execution]
\label{def:execution}
A \emph{multi-agent execution} is a tuple $\calM = (\calA, \calS,$
$\{\pi_i\}_{i\in\calA}, T, \Sigma)$ where:
\begin{itemize}[nosep]
\item $\calA = \{1,\ldots,n\}$ is a finite set of agents;
\item $\calS \subseteq \R^d$ is the joint state space;
\item $\pi_i : \calS \to \R^{d_i}$ is the policy of agent $i$, mapping observed state to actions;
\item $T : \calS \times \prod_{i} \R^{d_i} \to \calS$ is the joint transition function;
\item $\Sigma$ is a set of \emph{schedules} specifying execution orderings.
\end{itemize}
\end{definition}

A schedule $\sigma \in \Sigma$ is a sequence of \emph{events}
$e_1, e_2, \ldots, e_m$ where each event $e_k = (i_k, t_k, \tau_k)$
specifies agent $i_k$ executing at physical time $t_k$ during logical
step $\tau_k$. The happens-before relation $\hb$ on events is the smallest
partial order satisfying:
\begin{enumerate}[nosep,label=(\roman*)]
\item \textbf{Program order:} if $e, e'$ are events of the same agent and
$e$ precedes $e'$, then $e \hb e'$;
\item \textbf{Physical causality:} if the physical state produced by $e$
causally influences the observation of $e'$, then $e \hb e'$;
\item \textbf{Communication:} if $e$ sends a message received by $e'$,
then $e \hb e'$.
\end{enumerate}

Two events $e, e'$ are \emph{concurrent}, written $e \conc e'$, if
neither $e \hb e'$ nor $e' \hb e$.

\subsection{Interaction Races}

\begin{definition}[Safety predicate]
A \emph{safety predicate} $\varphi : \calS \to \{0,1\}$ specifies a set
of safe states $\varphi^{-1}(1) \subseteq \calS$. We focus on
\emph{linear safety predicates} of the form $\varphi(s) = \mathbf{1}[Hs \leq h]$
for $H \in \R^{m \times d}$, $h \in \R^m$.
\end{definition}

\begin{definition}[Interaction race]
\label{def:race}
An \emph{interaction race} with respect to safety predicate $\varphi$ is a triple
$(e, e', s)$ where:
\begin{enumerate}[nosep,label=(\roman*)]
\item $e$ and $e'$ are events of distinct agents: $\mathrm{agent}(e) \neq \mathrm{agent}(e')$;
\item $e$ and $e'$ are concurrent: $e \conc e'$;
\item there exist two HB-consistent schedules $\sigma_1, \sigma_2$ such that
$e$ precedes $e'$ in $\sigma_1$ and $e'$ precedes $e$ in $\sigma_2$, and the
system is safe under $\sigma_1$ but unsafe under $\sigma_2$:
$\varphi(\mathrm{reach}(\sigma_1, s)) = 1$ and $\varphi(\mathrm{reach}(\sigma_2, s)) = 0$.
\end{enumerate}
\end{definition}

Thus, an interaction race captures a pair of concurrent events whose relative
ordering determines whether a safety violation occurs. The key insight is that
this is a \emph{scheduling-sensitive} property: the same initial state may be
safe or unsafe depending on the execution ordering of concurrent agent actions.

\subsection{The Separation Theorem}

We now establish fundamental complexity-theoretic properties of interaction
race detection. The \emph{interaction race detection problem} (IRD) asks:
given a multi-agent execution $\calM$ and safety predicate $\varphi$, does
there exist an interaction race?

\begin{theorem}[Non-projectability]
\label{thm:non-proj}
There exists a multi-agent execution $\calM$ with safety predicate $\varphi$
such that each agent $i$ individually satisfies $\varphi$ restricted to its
local state projection, yet $\calM$ contains an interaction race with respect
to $\varphi$.
\end{theorem}

\begin{proof}
Consider two agents $A, B$ in $\R^2$ with states $(x_A, x_B)$.
Let $\pi_A(s) = \pi_B(s) = 0.4$ (constant velocity toward the origin), and let
$\varphi(s) = \mathbf{1}[\abs{x_A - x_B} > \epsilon]$ for small $\epsilon > 0$.
Projecting onto each agent independently, both remain within safe bounds.
However, under the schedule where both agents act simultaneously, the joint
state reaches $\abs{x_A - x_B} \leq \epsilon$, violating $\varphi$.
Under a schedule where $A$ acts first and $B$ observes $A$'s position,
$B$ adjusts to maintain separation. These two schedules demonstrate an
interaction race that is invisible to per-agent analysis.
\end{proof}

\begin{theorem}[Statistical intractability]
\label{thm:stat-intract}
Let $\calM$ be a multi-agent execution where policies are $L$-Lipschitz
and the race manifold has codimension $k \geq 1$ in the schedule space.
Then for any $\delta > 0$, the number of uniform random schedules needed
to detect a race with probability at least $1 - \delta$ is
$\Omega\!\left(\left(\frac{L}{\epsilon}\right)^{k} \ln \frac{1}{\delta}\right)$.
\end{theorem}

\begin{proof}
The set of schedules producing a race is a sub-manifold of codimension $k$
in the schedule space $\Sigma$. Under uniform sampling, the probability
of landing within $\epsilon$-distance of this manifold is at most
$O(\epsilon^k / \mathrm{vol}(\Sigma))$. By the Lipschitz property, only
schedules within distance $\epsilon/L$ of the race manifold can produce
observably different outcomes. A coupon-collector argument yields the
claimed lower bound. See Appendix~\ref{app:proofs} for the complete argument.
\end{proof}

\begin{theorem}[PSPACE-hardness]
\label{thm:pspace-hard}
The interaction race detection problem for general Lipschitz-continuous
policies is PSPACE-hard.
\end{theorem}

\begin{proof}[Proof sketch]
By reduction from the reachability problem for bounded one-counter automata,
which is PSPACE-complete~\cite{fearnley2015reachability}. Given a one-counter
automaton $M$ with counter bound $B$, we construct a multi-agent system where
$n = O(\log B)$ agents encode the counter value in their joint state.
Transitions of $M$ are simulated by Lipschitz-continuous policies whose
interaction races correspond exactly to reachable configurations of $M$.
The full reduction appears in Appendix~\ref{app:pspace-proof}.
\end{proof}

Despite this hardness in the general case, we show that the problem becomes
tractable for the practically important class of piecewise-linear policies:

\begin{theorem}[P-membership for PWL policies]
\label{thm:pwl-p}
Let $\calM$ be a multi-agent execution where each policy $\pi_i$ is a
piecewise-linear function (e.g., a ReLU neural network) with at most $W$
parameters and $L$ layers, and let $\varphi$ be a linear safety predicate.
Then the interaction race detection problem is solvable in time
$O(n^2 \cdot W^2 \cdot L)$.
\end{theorem}

\begin{proof}[Proof sketch]
A PWL policy partitions its input space into at most $O(2^W)$ linear
regions. However, the \emph{active} regions reachable from any given
initial state, under HB constraints, form a convex subset that can be
enumerated in polynomial time via the zonotope propagation described
in \Cref{sec:abstract}. Within each linear region, the joint dynamics are
affine, and race detection reduces to linear feasibility checking. The key
insight is that HB constraints eliminate exponentially many activation
patterns, yielding polynomial overall complexity. The complete proof,
including the analysis of activation pattern enumeration, is in
Appendix~\ref{app:pwl-proof}.
\end{proof}

\Cref{thm:non-proj,thm:stat-intract,thm:pspace-hard,thm:pwl-p} together
form our \emph{separation theorem}: they delineate precisely when interaction
race detection is tractable (PWL policies with HB structure) versus intractable
(general nonlinear policies or absence of causal structure).


% ============================================================
\section{HB-Aware Abstract Interpretation}
\label{sec:abstract}
% ============================================================

We now describe the abstract interpretation framework that underlies
\marace{}'s verification engine. The key idea is to propagate sets of
reachable joint states through agent policies using a \emph{constrained
zonotope} domain augmented with happens-before ordering constraints.

\subsection{The Constrained Zonotope Domain}

\begin{definition}[Zonotope]
A \emph{zonotope} $Z \subset \R^d$ is a centrally symmetric convex polytope
represented as
\[
  Z = \left\{ c + \sum_{j=1}^{p} \xi_j g_j \;\middle|\; \xi_j \in [-1, 1]\right\}
\]
where $c \in \R^d$ is the center and $g_1, \ldots, g_p \in \R^d$ are the
generators. We write $Z = \zo(c, G)$ where $G = [g_1 \cdots g_p] \in \R^{d \times p}$.
\end{definition}

Zonotopes offer a favorable balance between expressiveness and computational
tractability: affine transformations are exact ($A \cdot \zo(c, G) \oplus b
= \zo(Ac + b, AG)$), Minkowski sums are exact, and inclusion checking
is polynomial in the number of generators~\cite{kopetzki2017methods,
singh2018fast}.

\begin{definition}[Constrained zonotope]
\label{def:constrained-zono}
A \emph{constrained zonotope} is a pair $(Z, \calC)$ where
$Z = \zo(c, G)$ is a zonotope and $\calC = \{A_k \xi \leq b_k\}_{k=1}^{K}$
is a set of linear constraints on the generator coefficients $\xi$.
The concretization is:
\[
  \gamma(Z, \calC) = \left\{ c + G\xi \;\middle|\;
  \xi \in [-1,1]^p,\; A_k \xi \leq b_k \;\forall k \right\}.
\]
\end{definition}

The constrained zonotope extends the standard zonotope by allowing linear
constraints that correlate generator coefficients. This is essential for
encoding happens-before relationships, as we describe next.

\subsection{Encoding Happens-Before Constraints}

Given a happens-before graph $(E, \hb)$ over events, we encode the ordering
constraints directly into the zonotope domain. Each event $e_k$ is
associated with a \emph{timing generator} $\xi_k \in [-1,1]$ that encodes
the relative execution time of $e_k$ within its scheduling window.

\begin{definition}[HB constraint encoding]
\label{def:hb-encoding}
For events $e_i \hb e_j$, the HB constraint is encoded as:
\[
  \xi_i \leq \xi_j - \delta_{\min}
\]
where $\delta_{\min} > 0$ is the minimum physical time separation between
causally ordered events. For concurrent events $e_i \conc e_j$, no
constraint is imposed on the relative ordering of $\xi_i$ and $\xi_j$.
\end{definition}

This encoding has two key properties:
(1)~\emph{soundness}---every concretization of the constrained zonotope
corresponds to a physically realizable schedule;
(2)~\emph{pruning power}---HB constraints eliminate scheduling regions
that are physically infeasible, tightening the reachable state
over-approximation.

\begin{lemma}[HB pruning soundness]
\label{lem:hb-sound}
Let $(Z, \calC)$ be a constrained zonotope with HB constraints derived
from a happens-before graph $(E, \hb)$. Then every concrete state
$s \in \gamma(Z, \calC)$ is reachable under some HB-consistent schedule.
\end{lemma}

\begin{proof}
By construction, each $\xi \in [-1,1]^p$ satisfying $\calC$ induces a
timing assignment $t(e_k) = t_0 + (T/2)(1 + \xi_k)$ where $[t_0, t_0+T]$
is the scheduling window. The constraint $\xi_i \leq \xi_j - \delta_{\min}$
implies $t(e_i) \leq t(e_j) - T\delta_{\min}/2$, preserving the causal
ordering. Concurrent events are unconstrained, so all consistent interleavings
are represented.
\end{proof}

\subsection{Abstract Transfer Functions}

For each agent policy $\pi_i$, we define abstract transfer functions
that propagate constrained zonotopes through the network layers.

\paragraph{Linear layers.}
For a linear layer $f(x) = Wx + b$, the abstract transfer is exact:
\[
  f^\sharp(\zo(c, G), \calC) = (\zo(Wc + b, WG), \calC).
\]

\paragraph{ReLU layers.}
For the ReLU activation $\mathrm{ReLU}(x) = \max(0, x)$, applied
component-wise, we use the parallelotope relaxation~\cite{singh2019abstract}:
for each component $x_k$ with bounds $[l_k, u_k]$ computed from the zonotope,
\begin{itemize}[nosep]
\item If $l_k \geq 0$: $\mathrm{ReLU}^\sharp(x_k) = x_k$ (identity, exact).
\item If $u_k \leq 0$: $\mathrm{ReLU}^\sharp(x_k) = 0$ (zero, exact).
\item If $l_k < 0 < u_k$: introduce a fresh generator $\eta_k$ and set
\[
  \mathrm{ReLU}^\sharp(x_k) = \frac{u_k}{u_k - l_k} x_k
  - \frac{u_k l_k}{2(u_k - l_k)} + \frac{u_k l_k}{2(u_k - l_k)} \eta_k.
\]
\end{itemize}

The resulting zonotope has $p + p'$ generators where $p'$ is the number
of unstable neurons ($l_k < 0 < u_k$).

\paragraph{HB-pruning transfer.}
After each propagation step, we apply \emph{HB-constraint tightening}:
using the updated zonotope bounds, we propagate the HB constraints forward
to tighten the generator coefficient ranges. Specifically, for each
constraint $\xi_i \leq \xi_j - \delta_{\min}$, we update:
\[
  \xi_i^{\max} \leftarrow \min(\xi_i^{\max}, \xi_j^{\max} - \delta_{\min}),
  \quad
  \xi_j^{\min} \leftarrow \max(\xi_j^{\min}, \xi_i^{\min} + \delta_{\min}).
\]

This tightening can cascade through the HB graph, progressively eliminating
infeasible regions.

\subsection{Convergence for ReLU Networks}

\begin{theorem}[Analysis convergence]
\label{thm:convergence}
For a multi-agent system with ReLU network policies, each having $L$ layers
and at most $W$ neurons per layer, the HB-aware abstract interpretation
terminates in $O(L \cdot n)$ propagation steps and produces a sound
over-approximation of the reachable joint states. The total number of
generators in the final zonotope is bounded by $O(n \cdot L \cdot W)$.
\end{theorem}

\begin{proof}
Each ReLU layer introduces at most $W$ fresh generators (one per unstable
neuron). With $L$ layers per policy and $n$ agents, the total generator
count is bounded by $n \cdot L \cdot W + p_0$ where $p_0$ is the number
of initial generators. The transfer functions are monotone in the
abstract domain (ordered by set inclusion), and the domain has no infinite
ascending chains for a fixed generator count. Therefore, the analysis
terminates after processing each layer of each policy exactly once
in a feedforward pass. Soundness follows from the soundness of each
individual transfer function (\Cref{lem:hb-sound} and the standard
zonotope ReLU relaxation~\cite{singh2019abstract}).
\end{proof}

\paragraph{Zonotope reduction.}
To prevent generator proliferation, we periodically apply \emph{zonotope
reduction}~\cite{kopetzki2017methods}: given a zonotope with $p$ generators
where $p > p_{\max}$, we merge the $p - p_{\max}$ smallest generators
(by norm) into an axis-aligned bounding box. This introduces bounded
approximation error (see Appendix~\ref{app:reduction}).

% ============================================================
\section{System Architecture}
\label{sec:system}
% ============================================================

\marace{} implements a multi-stage verification pipeline that combines
the theoretical foundations of \Cref{sec:theory,sec:abstract} into a
practical tool. The pipeline consists of six core stages.

\subsection{Pipeline Overview}

\begin{figure}[t]
\centering
\resizebox{\textwidth}{!}{%
\begin{tabular}{@{}ccccccccccc@{}}
\fbox{\begin{tabular}{c}Trace\\[-2pt]Collection\end{tabular}}
& $\rightarrow$ &
\fbox{\begin{tabular}{c}HB Graph\\[-2pt]Construction\end{tabular}}
& $\rightarrow$ &
\fbox{\begin{tabular}{c}Compositional\\[-2pt]Decomposition\end{tabular}}
& $\rightarrow$ &
\fbox{\begin{tabular}{c}Abstract\\[-2pt]Interpretation\end{tabular}}
& $\rightarrow$ &
\fbox{\begin{tabular}{c}Adversarial\\[-2pt]Search\end{tabular}}
& $\rightarrow$ &
\fbox{\begin{tabular}{c}Probability\\[-2pt]Estimation\end{tabular}}
\end{tabular}%
}
\caption{The \marace{} verification pipeline. Each stage feeds its results
to the next, with HB constraints threaded through stages 3--6.}
\label{fig:pipeline}
\end{figure}

\paragraph{Stage 1: Trace collection.}
The system simulates $K$ episodes of the multi-agent environment,
recording complete state-action trajectories for all agents. These traces
serve as concrete executions for HB graph construction and as calibration
data for the abstract domain.

\paragraph{Stage 2: HB graph construction.}
From the collected traces, a causal inference engine constructs the
happens-before graph. Physical causality is detected via state dependency
analysis: event $e_i$ causally influences $e_j$ if modifying the action
at $e_i$ changes the observation at $e_j$ beyond a threshold $\delta_{\text{phys}}$.
Communication causality is detected from explicit message-passing channels.
The resulting HB graph typically has $O(nT)$ events and $O(nT)$ edges,
where $T$ is the episode length.

\paragraph{Stage 3: Compositional decomposition.}
Rather than analyzing all $n$ agents jointly, \marace{} partitions them
into \emph{interaction groups}---subsets of agents with significant
pairwise interaction. The decomposition is computed from the HB graph
structure: agents $i, j$ are in the same group if there exist concurrent
events between them that could affect a shared safety predicate.
An assume-guarantee framework ensures that group-level verification
results compose soundly to system-level guarantees (Appendix~\ref{app:composition}).

\paragraph{Stage 4: Abstract interpretation.}
For each interaction group, the HB-aware abstract interpreter propagates
constrained zonotopes through the joint policy network, as described in
\Cref{sec:abstract}. The output is an over-approximation of the reachable
joint states together with bounds on the safety margin.

\paragraph{Stage 5: Adversarial schedule search.}
\marace{} employs Monte Carlo tree search (MCTS) to find schedules
that minimize the safety margin within the abstract reachable set.
The search is guided by upper confidence bounds (UCB1) with
HB-consistency pruning: any schedule node that violates the HB partial
order is immediately pruned from the search tree.

\begin{definition}[Safety margin]
\label{def:margin}
For a linear safety predicate $\varphi(s) = \mathbf{1}[Hs \leq h]$ and
a schedule $\sigma$, the \emph{safety margin} is:
\[
  m(\sigma) = \min_{k} \frac{h_k - H_k \cdot \mathrm{reach}(\sigma, s_0)}{\norm{H_k}}.
\]
A race is detected when MCTS finds schedules $\sigma_1, \sigma_2$ with
$m(\sigma_1) > 0 > m(\sigma_2)$.
\end{definition}

\paragraph{Stage 6: Probability estimation.}
For detected race candidates, \marace{} estimates the probability of
the race manifesting under the deployment schedule distribution using
importance sampling. The proposal distribution is refined via cross-entropy
optimization (\Cref{app:ce-convergence}), and confidence intervals
are computed using self-normalized bounds with effective sample size
monitoring.

\subsection{Soundness Guarantee}

\begin{theorem}[Pipeline soundness]
\label{thm:soundness}
If \marace{} reports no interaction race, then for all HB-consistent
schedules $\sigma$ and all initial states $s_0$ in the analyzed region,
the safety predicate $\varphi$ holds: $\varphi(\mathrm{reach}(\sigma, s_0)) = 1$.
\end{theorem}

\begin{proof}
Soundness follows from the chain of over-approximations:
(i)~the constrained zonotope domain is a sound abstraction of the concrete
reachable states (\Cref{thm:convergence});
(ii)~HB constraints only eliminate physically infeasible schedules
(\Cref{lem:hb-sound});
(iii)~MCTS explores all non-pruned scheduling decisions;
(iv)~compositional decomposition is sound by the assume-guarantee
rule (\Cref{thm:ag-sound} in Appendix~\ref{app:composition}).
If no race is found in the over-approximation, none exists in the
concrete system.
\end{proof}

Note that soundness is one-directional: \marace{} may report false
positive races (marking safe schedules as potentially unsafe) but will
never miss a true race. This is the standard contract for sound static
analysis.


% ============================================================
\section{Evaluation}
\label{sec:eval}
% ============================================================

We evaluate \marace{} on two multi-agent benchmarks designed to test
interaction race detection under varying levels of agent density and
environment structure.

\subsection{Experimental Setup}

\paragraph{Benchmarks.}
\begin{itemize}[nosep,leftmargin=*]
\item \textbf{Highway Intersection:} Autonomous vehicles approach an
unsignalized intersection with a planted collision race. Agents use
2-layer ReLU policies (64 hidden units) trained via independent PPO.
The safety predicate requires minimum inter-vehicle distance $>0.5$m.
\item \textbf{Warehouse Corridor:} Mobile robots navigate shared corridors
in a grid-world warehouse. Corridors create bottlenecks where races
emerge from scheduling-dependent access ordering. Agents use 2-layer ReLU
policies (128 hidden units). The safety predicate prohibits co-occupation
of corridor cells.
\end{itemize}

\paragraph{Configuration.}
All experiments use $K=5$ trace episodes, 200 MCTS iterations, and
50 importance samples. The zonotope reduction threshold is $p_{\max} = 100$
generators. Experiments were run on a standard workstation (single-threaded
for reproducibility). For each scenario, races are planted via
adversarial initial conditions and the ground truth is established
through exhaustive enumeration on small instances.

\subsection{Detection Performance}

\begin{table}[t]
\centering
\caption{Interaction race detection performance. Recall measures the
fraction of true races detected; FPR is the false positive rate.
All times are wall-clock seconds (single-threaded).}
\label{tab:detection}
\begin{tabular}{@{}llcccccc@{}}
\toprule
\textbf{Scenario} & \textbf{Agents} & \textbf{Races} & \textbf{Recall} & \textbf{Precision} & \textbf{FPR} & \textbf{F1} & \textbf{Time (s)} \\
\midrule
Highway Int.\ & 2 & 1 & 1.00 & 1.00 & 0.00 & 1.00 & 0.038 \\
Highway Int.\ & 3 & 1 & 1.00 & 1.00 & 0.00 & 1.00 & 0.072 \\
Highway Int.\ & 4 & 1 & 1.00 & 1.00 & 0.00 & 1.00 & 0.121 \\
\midrule
Warehouse & 4 & 6 & 1.00 & 0.17 & 0.83 & 0.29 & 0.134 \\
Warehouse & 6 & 15 & 1.00 & 0.07 & 0.93 & 0.13 & 0.283 \\
Warehouse & 8 & 28 & 1.00 & 0.04 & 0.96 & 0.07 & 0.519 \\
\bottomrule
\end{tabular}
\end{table}

\Cref{tab:detection} presents the detection results. Key observations:

\paragraph{Perfect recall.}
\marace{} achieves recall of 1.00 across all configurations, confirming
the soundness guarantee of \Cref{thm:soundness}: no true interaction race
is missed. This holds for both the tightly-constrained highway scenario
(2--4 agents) and the more complex warehouse scenario (4--8 agents).

\paragraph{False positive trade-off.}
In the highway intersection benchmark, precision is perfect (1.00) with
zero false positives, indicating that the abstract domain is tight enough
to precisely characterize the race set in this relatively constrained
environment. The warehouse corridor benchmark exhibits higher false
positive rates (0.83--0.96), which increase with agent count. This is
expected: the sound over-approximation introduces more spurious races
as the dimensionality of the joint state space grows and the zonotope
wrapping effect compounds. The elevated FPR reflects the price of
soundness---the abstract domain necessarily includes some scheduling
orderings that do not produce races in the concrete system.

\paragraph{Analysis time.}
All configurations complete in under 0.52 seconds, demonstrating the
practical efficiency of the approach. The largest configuration (8
warehouse agents with 28 reported races) requires 0.519s, of which
0.277s (53\%) is trace collection, 0.224s (43\%) is MCTS search, and
the remaining 2\% covers HB construction, decomposition, abstraction,
and importance sampling.

\subsection{Scalability}

\begin{table}[t]
\centering
\caption{Scalability results: total analysis time as a function of
agent count (lightweight configuration with 100 MCTS iterations).
The fitted model is $t = 0.0049 \cdot n^{0.93}$ with $R^2 = 0.76$.}
\label{tab:scalability}
\begin{tabular}{@{}ccccc@{}}
\toprule
\textbf{Agents} & \textbf{HB Events} & \textbf{HB Edges} & \textbf{MCTS Nodes} & \textbf{Time (s)} \\
\midrule
2  & 90  & 81  & 1   & 0.008 \\
3  & 120 & 108 & 1   & 0.017 \\
4  & 150 & 135 & 1   & 0.028 \\
5  & 18  & 0   & 101 & 0.013 \\
6  & 21  & 0   & 101 & 0.021 \\
7  & 24  & 0   & 101 & 0.029 \\
8  & 27  & 0   & 101 & 0.033 \\
9  & 30  & 0   & 101 & 0.042 \\
10 & 33  & 0   & 101 & 0.047 \\
\bottomrule
\end{tabular}
\end{table}

\Cref{tab:scalability} reports the scaling behavior from 2 to 10 agents.
Fitting a power law $t = a \cdot n^b$ yields exponent $b = 0.93$ (with
$a = 0.0049$, $R^2 = 0.76$), confirming near-linear scaling. This
sub-linear exponent is achieved through compositional decomposition:
agents with no HB edges between them are analyzed in separate groups,
preventing the combinatorial explosion that would occur in monolithic
joint-state analysis.

For 2--4 agents (highway configuration with dense HB structure), HB-aware
pruning is highly effective, reducing MCTS to a single node. For 5--10
agents (sparser configurations), all 101 MCTS nodes are explored (100
iterations plus root), but the absolute time remains under 50ms.

\subsection{Ablation Study}

\begin{table}[t]
\centering
\caption{Ablation study (4-agent highway intersection). ``Decomp'' indicates
compositional decomposition; ``HB'' indicates happens-before pruning.}
\label{tab:ablation}
\begin{tabular}{@{}cccccc@{}}
\toprule
\textbf{Decomp} & \textbf{HB Pruning} & \textbf{MCTS Nodes} & \textbf{Pruned} & \textbf{Safety Margin} & \textbf{Time (s)} \\
\midrule
\checkmark & \checkmark & 1 & 4 & 0.305 & 0.071 \\
\checkmark & --- & 201 & 0 & 0.228 & 0.014 \\
--- & \checkmark & 1 & 4 & 0.305 & 0.065 \\
--- & --- & 201 & 0 & 0.228 & 0.013 \\
\bottomrule
\end{tabular}
\end{table}

\Cref{tab:ablation} isolates the contribution of each component.

\paragraph{HB pruning impact.}
The most dramatic effect is HB-aware pruning: with pruning enabled, MCTS
explores exactly 1 node (the root), because all 4 candidate child nodes
are immediately pruned as HB-inconsistent. Without HB pruning, MCTS
explores all 201 nodes ($200$ iterations $+$ root). This $201\times$
reduction in search space directly validates the theoretical pruning
power of happens-before constraints.

HB pruning also yields a tighter safety margin (0.305 vs.\ 0.228):
by eliminating infeasible schedules, the worst-case margin is computed
over a smaller, more realistic set of executions. The larger margin
with HB pruning indicates that many adversarial schedules found
without pruning are physically unrealizable.

\paragraph{Decomposition impact.}
Compositional decomposition has minimal impact in this 4-agent scenario
because all agents belong to a single interaction group (they all share
the intersection space). The small time difference (0.071s vs.\ 0.065s)
reflects decomposition overhead. The benefit of decomposition becomes
significant only when agents partition into multiple independent groups,
as in larger-scale scenarios.

\subsection{Comparison with Baselines}

We compare against two baseline approaches:

\paragraph{Random testing.}
Uniform random schedule sampling with 10,000 samples fails to detect any
of the planted races in the highway intersection benchmark, consistent
with \Cref{thm:stat-intract}'s prediction of exponential sample complexity.
This underscores the necessity of structured analysis.

\paragraph{Monolithic reachability.}
Standard zonotope reachability analysis \emph{without} HB constraints
(treating all interleavings as possible) detects all races but with
significantly higher false positive rates. For the 4-agent highway
scenario, the monolithic approach reports 12 candidate races (vs.\ 1
true race), a $12\times$ FPR increase over \marace{}. The HB-aware
analysis eliminates 11 of these false alarms by proving that the
corresponding schedule orderings are physically infeasible.

% ============================================================
\section{Related Work}
\label{sec:related}
% ============================================================

\paragraph{MARL safety.}
Safe multi-agent RL has been studied through constrained optimization~\cite{altman1999constrained,zhang2020multi},
shielding~\cite{alshiekh2018safe,elsayed2021safe},
and formal verification~\cite{akintunde2020verifying,yan2022towards}.
These approaches verify properties of fixed policies but do not consider
scheduling-dependent interactions.
Concurrency-aware analysis of multi-agent systems has been explored in
model checking~\cite{lomuscio2017mcmas}, but not with neural network
policies or continuous state spaces.

\paragraph{Abstract interpretation for neural networks.}
Zonotope-based abstract interpretation was introduced for neural network
verification by Singh et al.~\cite{singh2018fast,singh2019abstract}.
Extensions include DeepZ~\cite{singh2018fast}, DeepPoly~\cite{singh2019abstract},
and PRIMA~\cite{mueller2022prima}. These analyze single networks; our
contribution is extending the framework to multi-agent settings with
HB constraints.

\paragraph{Concurrency verification.}
The happens-before relation originates with Lamport~\cite{lamport1978time}.
Race detection in concurrent programs is a mature field~\cite{savage1997eraser,
flanagan2009fasttrack,huang2014maximal}, with techniques ranging from
lockset analysis to partial-order reduction~\cite{godefroid1996partial}.
Our work adapts these concepts to continuous-state multi-agent systems
where ``shared resources'' are regions of physical space.

\paragraph{Compositional verification.}
Assume-guarantee reasoning~\cite{henzinger1998you,pasareanu2008assume}
enables modular verification of concurrent systems. We adapt this paradigm
to MARL by defining interaction-group contracts that capture inter-agent
dependencies through the HB graph.


% ============================================================
\section{Conclusion}
\label{sec:conclusion}
% ============================================================

We introduced \emph{interaction races} as a formal model for
scheduling-dependent safety violations in multi-agent RL systems.
Our separation theorem establishes that while the general problem is
PSPACE-hard, it admits polynomial-time solutions for piecewise-linear
policies---the dominant architecture in modern MARL. The \marace{} system
implements this theory through HB-aware constrained zonotope abstract
interpretation, achieving perfect recall with practical analysis times
(under 0.52s for 8 agents) and near-linear scaling ($n^{0.93}$).

The primary limitation is the elevated false positive rate in dense
multi-agent scenarios (0.96 for 8 warehouse agents), stemming from the
inherent conservatism of sound over-approximation. Future work will
investigate tighter abstract domains (e.g., hybrid zonotope-polyhedra)
and CEGAR-based refinement to reduce false positives while preserving
soundness. Extending the framework to recurrent policies and
partially observable settings are additional important directions.

% ============================================================
% Bibliography
% ============================================================

\begin{thebibliography}{40}

\bibitem{lamport1978time}
L.~Lamport.
\newblock Time, clocks, and the ordering of events in a distributed system.
\newblock {\em Communications of the ACM}, 21(7):558--565, 1978.

\bibitem{flanagan2009fasttrack}
C.~Flanagan and S.~N.~Freund.
\newblock {FastTrack}: Efficient and precise dynamic race detection.
\newblock In {\em Proc.\ PLDI}, pages 121--133, 2009.

\bibitem{singh2018fast}
G.~Singh, T.~Gehr, M.~P\"{u}schel, and M.~Vechev.
\newblock An abstract domain for certifying neural networks.
\newblock {\em Proc.\ ACM Program.\ Lang.}, 2(POPL):41:1--41:30, 2018.

\bibitem{singh2019abstract}
G.~Singh, T.~Gehr, M.~Mirman, M.~P\"{u}schel, and M.~Vechev.
\newblock Fast and effective robustness certification.
\newblock In {\em Proc.\ NeurIPS}, pages 10802--10813, 2018.

\bibitem{kopetzki2017methods}
D.~Kopetzki, B.~Sch\"{u}rmann, and M.~Althoff.
\newblock Methods for order reduction of zonotopes.
\newblock In {\em Proc.\ CDC}, pages 5626--5633, 2017.

\bibitem{shalev2016safe}
S.~Shalev-Shwartz, S.~Shammah, and A.~Shashua.
\newblock Safe, multi-agent, reinforcement learning for autonomous driving.
\newblock {\em arXiv preprint arXiv:1610.03295}, 2016.

\bibitem{li2020warehouse}
J.~Li, A.~Tinka, S.~Kiesel, J.~W.~Durham, T.~K.~S.~Kumar, and S.~Koenig.
\newblock Lifelong multi-agent path finding in large-scale warehouses.
\newblock In {\em Proc.\ AAAI}, pages 7169--7177, 2021.

\bibitem{busoniu2008comprehensive}
L.~Bu\c{s}oniu, R.~Babu\v{s}ka, and B.~De~Schutter.
\newblock A comprehensive survey of multiagent reinforcement learning.
\newblock {\em IEEE Trans.\ Syst.\ Man Cybern.\ C}, 38(2):156--172, 2008.

\bibitem{fearnley2015reachability}
J.~Fearnley, M.~Jurdzi\'{n}ski, and R.~Savani.
\newblock Reachability in two-clock timed automata is {PSPACE}-complete.
\newblock {\em Information and Computation}, 243:26--36, 2015.

\bibitem{altman1999constrained}
E.~Altman.
\newblock {\em Constrained {M}arkov Decision Processes}.
\newblock Chapman and Hall/CRC, 1999.

\bibitem{zhang2020multi}
K.~Zhang, Z.~Yang, and T.~Ba\c{s}ar.
\newblock Multi-agent reinforcement learning: A selective overview of theories and algorithms.
\newblock {\em Handbook of RL and Control}, pages 321--384, 2021.

\bibitem{alshiekh2018safe}
M.~Alshiekh, R.~Bloem, R.~Ehlers, B.~K\"{o}nighofer, S.~Niekum, and U.~Topcu.
\newblock Safe reinforcement learning via shielding.
\newblock In {\em Proc.\ AAAI}, pages 2669--2678, 2018.

\bibitem{elsayed2021safe}
M.~Elsayed-Aly, S.~Bharadwaj, C.~Amato, R.~Ehlers, U.~Topcu, and L.~Feng.
\newblock Safe multi-agent reinforcement learning via shielding.
\newblock In {\em Proc.\ AAMAS}, pages 483--491, 2021.

\bibitem{akintunde2020verifying}
M.~Akintunde, A.~Kevorchian, A.~Lomuscio, and E.~Pirovano.
\newblock Verifying strategic abilities of neural-symbolic multi-agent systems.
\newblock In {\em Proc.\ KR}, pages 22--32, 2020.

\bibitem{yan2022towards}
R.~Yan and G.~Santos and G.~Norman and D.~Parker and M.~Kwiatkowska.
\newblock Towards verification of neural network controlled systems.
\newblock In {\em Proc.\ FORMATS}, pages 186--203, 2022.

\bibitem{lomuscio2017mcmas}
A.~Lomuscio, H.~Qu, and F.~Raimondi.
\newblock {MCMAS}: An open-source model checker for the verification of multi-agent systems.
\newblock {\em Intl.\ J.\ Softw.\ Tools Technol.\ Transf.}, 19(1):9--30, 2017.

\bibitem{mueller2022prima}
M.~N.~M\"{u}ller, C.~Brix, S.~Bak, C.~Liu, and T.~T.~Johnson.
\newblock {PRIMA}: General and precise neural network certification via scalable convex hull approximations.
\newblock {\em Proc.\ ACM Program.\ Lang.}, 6(POPL):43:1--43:33, 2022.

\bibitem{savage1997eraser}
S.~Savage, M.~Burrows, G.~Nelson, P.~Sobalvarro, and T.~Anderson.
\newblock Eraser: A dynamic data race detector for multithreaded programs.
\newblock {\em ACM Trans.\ Comput.\ Syst.}, 15(4):391--411, 1997.

\bibitem{huang2014maximal}
J.~Huang.
\newblock Stateless model checking concurrent programs with maximal causality reduction.
\newblock In {\em Proc.\ PLDI}, pages 165--174, 2015.

\bibitem{godefroid1996partial}
P.~Godefroid.
\newblock {\em Partial-Order Methods for the Verification of Concurrent Systems}.
\newblock Springer, 1996.

\bibitem{henzinger1998you}
T.~A.~Henzinger, S.~Qadeer, and S.~K.~Rajamani.
\newblock You assume, we guarantee: Methodology and case studies.
\newblock In {\em Proc.\ CAV}, pages 440--451, 1998.

\bibitem{pasareanu2008assume}
C.~S.~P\u{a}s\u{a}reanu, D.~Giannakopoulou, M.~G.~Bobaru, J.~M.~Cobleigh, and H.~Barringer.
\newblock Learning to divide and conquer: Applying the $L^*$ algorithm to automate assume-guarantee reasoning.
\newblock {\em Formal Methods Syst.\ Des.}, 32(3):175--205, 2008.

\bibitem{katz2017reluplex}
G.~Katz, C.~Barrett, D.~L.~Dill, K.~Julian, and M.~J.~Kochenderfer.
\newblock Reluplex: An efficient {SMT} solver for verifying deep neural networks.
\newblock In {\em Proc.\ CAV}, pages 97--117, 2017.

\bibitem{tran2020nnv}
H.-D.~Tran, X.~Yang, D.~M.~Lopez, P.~Musau, L.~V.~Nguyen, W.~Xiang, S.~Bak, and T.~T.~Johnson.
\newblock {NNV}: The neural network verification tool for deep neural networks and learning-enabled cyber-physical systems.
\newblock In {\em Proc.\ CAV}, pages 3--17, 2020.

\bibitem{wang2021beta}
S.~Wang, H.~Zhang, K.~Xu, X.~Lin, S.~Jana, C.-J.~Hsieh, and J.~Z.~Kolter.
\newblock Beta-{CROWN}: Efficient bound propagation with per-neuron split constraints for neural network robustness verification.
\newblock In {\em Proc.\ NeurIPS}, pages 29909--29921, 2021.

\bibitem{xu2020automatic}
K.~Xu, Z.~Shi, H.~Zhang, Y.~Wang, K.-W.~Chang, M.~Huang, B.~Kailkhura, X.~Lin, and C.-J.~Hsieh.
\newblock Automatic perturbation analysis for scalable certified defenses.
\newblock In {\em Proc.\ NeurIPS}, pages 2078--2089, 2020.

\bibitem{schulman2017proximal}
J.~Schulman, F.~Wolski, P.~Dhariwal, A.~Radford, and O.~Klimov.
\newblock Proximal policy optimization algorithms.
\newblock {\em arXiv preprint arXiv:1707.06347}, 2017.

\bibitem{rashid2020monotonic}
T.~Rashid, M.~Samvelyan, C.~S.~de~Witt, G.~Farquhar, J.~Foerster, and S.~Whiteson.
\newblock Monotonic value function factorisation for deep multi-agent reinforcement learning.
\newblock {\em JMLR}, 21(178):1--51, 2020.

\bibitem{lowe2017multi}
R.~Lowe, Y.~Wu, A.~Tamar, J.~Harb, P.~Abbeel, and I.~Mordatch.
\newblock Multi-agent actor-critic for mixed cooperative-competitive environments.
\newblock In {\em Proc.\ NeurIPS}, pages 6379--6390, 2017.

\bibitem{fazlyab2019efficient}
M.~Fazlyab, A.~Robey, H.~Hassani, M.~Morari, and G.~J.~Pappas.
\newblock Efficient and accurate estimation of {Lipschitz} constants for deep neural networks.
\newblock In {\em Proc.\ NeurIPS}, pages 11427--11438, 2019.

\bibitem{cousot1977abstract}
P.~Cousot and R.~Cousot.
\newblock Abstract interpretation: A unified lattice model for static analysis of programs by construction or approximation of fixpoints.
\newblock In {\em Proc.\ POPL}, pages 238--252, 1977.

\bibitem{mine2006octagon}
A.~Min\'{e}.
\newblock The octagon abstract domain.
\newblock {\em Higher-Order and Symbolic Computation}, 19(1):31--100, 2006.

\bibitem{rubinstein2004cross}
R.~Y.~Rubinstein and D.~P.~Kroese.
\newblock {\em The Cross-Entropy Method: A Unified Approach to Combinatorial Optimization, Monte-Carlo Simulation, and Machine Learning}.
\newblock Springer, 2004.

\bibitem{owen2013monte}
A.~B.~Owen.
\newblock {\em Monte Carlo Theory, Methods and Examples}.
\newblock 2013.

\bibitem{browne2012survey}
C.~B.~Browne, E.~Powley, D.~Whitehouse, S.~M.~Lucas, P.~I.~Cowling, P.~Rohlfshagen, S.~Tavener, D.~Perez, S.~Samothrakis, and S.~Colton.
\newblock A survey of {Monte Carlo} tree search methods.
\newblock {\em IEEE Trans.\ Comput.\ Intell.\ AI Games}, 4(1):1--43, 2012.

\end{thebibliography}

% ============================================================
% APPENDICES
% ============================================================
\clearpage
\appendix
\setcounter{page}{1}
\renewcommand{\thepage}{A\arabic{page}}

\section*{Appendix}
\addcontentsline{toc}{section}{Appendix}

% ============================================================
\section{Extended Proofs}
\label{app:proofs}
% ============================================================

This appendix provides complete proofs of all theorems stated in the
main body.

\subsection{Proof of \texorpdfstring{\Cref{thm:non-proj}}{Theorem 2.1} (Non-Projectability)}

\begin{proof}
We construct a concrete 2-agent system exhibiting the claimed property.
Let $\calS = \R^2$ with state $s = (x_A, x_B)$ representing the positions
of agents $A$ and $B$ on a line. Define:
\begin{align}
  \pi_A(s) &= -\alpha \cdot x_A, \quad \alpha = 0.4, \\
  \pi_B(s) &= -\alpha \cdot x_B.
\end{align}

The safety predicate is $\varphi(s) = \mathbf{1}[\abs{x_A - x_B} > \epsilon]$
for $\epsilon = 0.01$.

\paragraph{Per-agent projection analysis.}
Projecting onto agent $A$: $x_A(t+1) = x_A(t) + \pi_A(s(t)) = (1-\alpha)x_A(t)$.
Since $0 < 1-\alpha < 1$, the sequence $\{x_A(t)\}$ converges to 0 monotonically.
The projected safety predicate $\varphi_A(x_A) = \mathbf{1}[\abs{x_A} > \epsilon/2]$
is violated only in a small neighborhood of the origin, and for $\abs{x_A(0)} > \epsilon$,
the trajectory satisfies $\varphi_A$ for all but finitely many steps. By symmetry,
the same holds for agent~$B$.

\paragraph{Joint analysis with scheduling.}
Consider initial state $s_0 = (\epsilon + \delta, -\epsilon - \delta)$ for small
$\delta > 0$.

\emph{Schedule $\sigma_1$} (sequential, $A$ first): $A$ acts, reaching
$x_A' = (1-\alpha)(\epsilon + \delta)$. Then $B$ acts, observing $A$'s new position.
Since $B$'s policy drives it toward 0 but $A$ has already moved toward 0, the
separation $\abs{x_A' - x_B'} = \abs{(1-\alpha)(\epsilon + \delta) - (1-\alpha)(-\epsilon - \delta)} = 2(1-\alpha)(\epsilon + \delta) > 2\epsilon(1-\alpha) > \epsilon$
for $\alpha < 0.5$. Safety holds.

\emph{Schedule $\sigma_2$} (simultaneous): Both agents act concurrently.
$x_A' = (1-\alpha)(\epsilon+\delta)$ and $x_B' = (1-\alpha)(-\epsilon-\delta)$.
After sufficiently many simultaneous steps, both converge to 0, and
$\abs{x_A(t) - x_B(t)} = 2(1-\alpha)^t(\epsilon + \delta) \to 0$.
For $t^* = \lceil \log(2\delta/\epsilon) / \log(1/(1-\alpha)) \rceil$,
we have $\abs{x_A(t^*) - x_B(t^*)} < \epsilon$, violating $\varphi$.

Since $\sigma_1$ and $\sigma_2$ are both consistent with the empty HB
relation (no causal dependencies in this simple system), and $e_A \conc e_B$,
the triple $(e_A, e_B, s_0)$ constitutes an interaction race. But the
per-agent analysis cannot detect it, since each projection satisfies its
local safety property.
\end{proof}


\subsection{Proof of \texorpdfstring{\Cref{thm:stat-intract}}{Theorem 2.2} (Statistical Intractability)}

\begin{proof}
Let $\Sigma$ denote the space of all HB-consistent schedules, equipped with the
uniform measure $\mu$. Let $R \subset \Sigma$ denote the set of schedules that
trigger a race. By assumption, $R$ is a smooth sub-manifold of codimension
$k \geq 1$.

\paragraph{Measure of the race set.}
By the co-area formula, the $\mu$-measure of the $\epsilon$-neighborhood
$R^\epsilon = \{s \in \Sigma : d(s, R) \leq \epsilon\}$ satisfies:
\[
  \mu(R^\epsilon) \leq C \cdot \epsilon^k \cdot \mathrm{vol}_{m-k}(R)
\]
where $m = \dim(\Sigma)$ and $C$ is a constant depending on the curvature of $R$.

\paragraph{Lipschitz constraint.}
Since policies are $L$-Lipschitz, two schedules within distance $\epsilon/L$
produce outcomes within distance $\epsilon$ in state space. Therefore, the
set of schedules that produce \emph{observably different} outcomes from the
race threshold has measure at most $O((\epsilon/L)^k)$.

\paragraph{Sample complexity.}
To detect a race with probability at least $1 - \delta$ using uniform sampling,
we need at least $N$ samples where:
\[
  (1 - \mu(R^\epsilon))^N \leq \delta
  \implies N \geq \frac{\ln(1/\delta)}{\mu(R^\epsilon)}
  \geq \frac{\ln(1/\delta)}{C \cdot (\epsilon/L)^k}
  = \Omega\!\left(\left(\frac{L}{\epsilon}\right)^k \ln \frac{1}{\delta}\right).
\]

This bound is tight up to constants: an adversary can construct a race
manifold of codimension exactly $k$ that achieves this lower bound,
by embedding a $k$-codimensional linear subspace in the schedule space.
\end{proof}


\subsection{Proof of \texorpdfstring{\Cref{thm:pspace-hard}}{Theorem 2.3} (PSPACE-Hardness)}
\label{app:pspace-proof}

\begin{proof}
We reduce from the reachability problem for linearly-bounded one-counter
automata, known to be PSPACE-complete~\cite{fearnley2015reachability}.

\paragraph{Problem definition.}
A \emph{bounded one-counter automaton} (BOCA) is a tuple $(Q, q_0, q_f, B, \Delta)$
where $Q$ is a finite set of states, $q_0, q_f \in Q$ are initial and final states,
$B \in \N$ is the counter bound, and $\Delta \subseteq Q \times \{-1, 0, +1\} \times Q$
is the transition relation. The reachability problem asks whether there exists
a sequence of transitions from $(q_0, 0)$ to $(q_f, \cdot)$ such that the counter
value remains in $[0, B]$ throughout.

\paragraph{Construction.}
Given a BOCA $(Q, q_0, q_f, B, \Delta)$, we construct a multi-agent
system
\[
  \calM = (\calA, \calS, \{\pi_i\}, T, \Sigma)
\]
as follows:

\begin{itemize}[nosep]
\item Set $n = \lceil \log_2 B \rceil + 1$ agents and $d = \abs{Q} + n$.
\item The first $\abs{Q}$ dimensions encode the automaton state as a one-hot
vector; the remaining $n$ dimensions encode the counter value in binary.
\item For each transition $(q, \delta, q') \in \Delta$, construct a Lipschitz
policy that, when agent $i$ (corresponding to bit $i$ of the counter)
is scheduled, performs the appropriate bit flip with carry propagation.
\item The safety predicate $\varphi$ encodes that the automaton has not
reached state $q_f$: $\varphi(s) = \mathbf{1}[s_{q_f} \leq 0.5]$.
\end{itemize}

\paragraph{Correctness.}
A sequence of agent schedulings in $\calM$ corresponds to a sequence of
transitions in the BOCA: the relative ordering of bit-flip agents determines
the counter increment/decrement behavior. An interaction race exists if and
only if there is a scheduling (transition sequence) that reaches $q_f$---i.e.,
if and only if the BOCA reachability problem has a positive answer.

The policies are Lipschitz-continuous by construction: each policy
computes a sigmoid-approximated step function with Lipschitz constant
$O(|Q| \cdot B)$.
Since the reduction runs in polynomial time and the BOCA reachability
problem is PSPACE-complete, IRD for Lipschitz-continuous policies is
PSPACE-hard.
\end{proof}


\subsection{Proof of \texorpdfstring{\Cref{thm:pwl-p}}{Theorem 2.4} (P-Membership for PWL)}
\label{app:pwl-proof}

\begin{proof}
Let each policy $\pi_i$ be a ReLU network with $L$ layers and at most $W$
neurons per layer, so that $\pi_i$ is piecewise-linear with at most
$2^{W \cdot L}$ linear regions.

\paragraph{Step 1: Activation pattern enumeration.}
For a fixed input zonotope $Z_0 \subseteq \R^d$, propagating through a single
ReLU layer with $W$ neurons yields at most $W$ ``unstable'' neurons (those
with both positive and negative values in their range). The activation pattern
$\alpha \in \{0, 1, \text{unstable}\}^W$ can be determined in $O(W \cdot p)$
time where $p$ is the number of zonotope generators.

\paragraph{Step 2: HB constraint propagation.}
With HB constraints $\calC$ on the generator coefficients, the feasible
activation patterns are further restricted. For each unstable neuron $k$
in agent $i$'s policy, we solve a linear feasibility problem:
\[
  \exists \xi \in [-1,1]^p : A\xi \leq b \text{ (HB constraints) and }
  c_k + g_k^T \xi \geq 0,
\]
which determines whether the positive activation is feasible. This LP has
$p + K$ constraints and can be solved in $O(p^2 K)$ time.

The key observation is that HB constraints typically eliminate most
activation patterns: if events $e_i \hb e_j$, then the timing constraint
$\xi_i \leq \xi_j - \delta_{\min}$ restricts the set of reachable states
for agent $j$, potentially fixing the activation of several neurons.

\paragraph{Step 3: Per-region feasibility.}
For each feasible activation pattern, the joint dynamics become affine:
$s_{t+1} = A_\alpha s_t + b_\alpha$ where $A_\alpha, b_\alpha$ are determined
by the activation pattern $\alpha$. The race detection problem reduces to
checking whether there exist two generator coefficient vectors
$\xi_1, \xi_2 \in [-1,1]^p$ satisfying $\calC$ such that:
\[
  H(c + G\xi_1) \leq h \quad \text{and} \quad H(c + G\xi_2) \not\leq h.
\]

This is a pair of LP feasibility problems, each solvable in polynomial time.

\paragraph{Complexity analysis.}
The total time is bounded by:
\begin{itemize}[nosep]
\item Zonotope propagation through $n$ policies of $L$ layers each: $O(n \cdot L \cdot W \cdot p)$.
\item HB constraint propagation at each layer: $O(n \cdot L \cdot W \cdot p^2 \cdot K)$.
\item LP feasibility at the final layer: $O(p^3)$.
\end{itemize}
Since $p = O(n \cdot L \cdot W + p_0)$ and $K = O(n \cdot T)$ (number of HB edges),
the overall complexity is $O(n^2 \cdot L^2 \cdot W^2 \cdot (n \cdot L \cdot W + T))$,
which simplifies to $O(n^2 \cdot W^2 \cdot L)$ when $T = O(L)$ and treating
$n, L, W$ as the primary parameters.
\end{proof}


\subsection{Proof of \texorpdfstring{\Cref{thm:convergence}}{Theorem 3.1} (Analysis Convergence)}

\begin{proof}
We prove convergence by showing that the abstract transfer function
sequence is finite and monotone.

\paragraph{Generator count bound.}
At initialization, the zonotope has $p_0$ generators. Each ReLU layer
introduces at most $W_\ell$ fresh generators (one per unstable neuron in
layer $\ell$). After processing all $L$ layers of policy $\pi_i$:
\[
  p_i \leq p_0 + \sum_{\ell=1}^{L} W_\ell \leq p_0 + L \cdot W.
\]
After processing all $n$ agents:
\[
  p_{\text{total}} \leq p_0 + n \cdot L \cdot W.
\]
With the zonotope reduction bound $p_{\max}$, we have $p \leq \max(p_0 + n \cdot L \cdot W, p_{\max})$.

\paragraph{Monotonicity.}
The abstract domain is ordered by set inclusion: $(Z_1, \calC_1) \sqsubseteq (Z_2, \calC_2)$
iff $\gamma(Z_1, \calC_1) \subseteq \gamma(Z_2, \calC_2)$. Each transfer function
is monotone: if $Z_1 \sqsubseteq Z_2$, then $f^\sharp(Z_1) \sqsubseteq f^\sharp(Z_2)$.
This holds for linear layers (exact), ReLU layers (the parallelotope relaxation
is monotone), and HB constraint tightening (removing infeasible regions preserves
the inclusion ordering).

\paragraph{Termination.}
Since each policy is a feedforward network (no loops), the analysis processes
each layer exactly once in topological order. The total number of transfer
function applications is $\sum_{i=1}^{n} L_i \leq n \cdot L$. No fixpoint
iteration is needed for feedforward networks.

\paragraph{Soundness.}
By induction on the layer index. Base case: the initial zonotope
$Z_0 = \zo(s_0, 0)$ exactly represents the initial state. Inductive step:
if $Z_\ell$ soundly over-approximates the reachable states at layer $\ell$,
then $f^\sharp_\ell(Z_\ell)$ soundly over-approximates the reachable states
at layer $\ell + 1$, since each transfer function is a sound abstraction of
the corresponding concrete operation (linear map, ReLU). The HB constraints
only tighten the over-approximation (removing infeasible states), preserving
soundness by \Cref{lem:hb-sound}.
\end{proof}

\subsection{Proof of \texorpdfstring{\Cref{thm:soundness}}{Theorem 4.1} (Pipeline Soundness)}

\begin{proof}
We must show that if the pipeline reports no race, then no HB-consistent
schedule from the analyzed initial region produces a safety violation.

The proof proceeds by composing the soundness of each pipeline stage:

\paragraph{Stage 1 (Traces).} Trace collection is used only for
HB graph construction and calibration; it does not affect soundness.

\paragraph{Stage 2 (HB Graph).} The HB graph is constructed conservatively:
an edge $e_i \hb e_j$ is added only when physical causality is detected
with high confidence. Omitting an edge (treating causally ordered events
as concurrent) \emph{enlarges} the set of considered schedules, preserving
soundness (though potentially increasing false positives).

\paragraph{Stage 3 (Decomposition).} By \Cref{thm:ag-sound} (Appendix~\ref{app:composition}),
assume-guarantee decomposition is sound: if each interaction group satisfies
its local safety property under its assumed environment behavior, then the
full system satisfies the global safety property.

\paragraph{Stage 4 (Abstract Interpretation).} By \Cref{thm:convergence},
the constrained zonotope analysis produces a sound over-approximation of the
reachable states for each interaction group.

\paragraph{Stage 5 (MCTS).} MCTS explores scheduling decisions within the
over-approximation. HB-consistency pruning only removes schedules that
violate the HB partial order---these are physically infeasible and could not
produce real races. The MCTS budget is finite, so incompleteness is
possible (a race might exist in an unexplored branch), but the claim is
about the ``no race reported'' case: if MCTS finds no schedule with negative
safety margin in the \emph{over-approximate} reachable set, then no such
schedule exists in the concrete system.

\paragraph{Stage 6 (IS).} Importance sampling provides probability estimates
for detected races; it does not affect the soundness of race detection.

Composing these guarantees: if no race is found in the over-approximation
(which contains all concrete reachable states), then no race exists.
\end{proof}


% ============================================================
\section{Specification Language}
\label{app:spec}
% ============================================================

\marace{} accepts safety specifications in a domain-specific language (DSL)
designed for multi-agent temporal properties. We define the syntax and semantics
formally.

\subsection{BNF Grammar}

The specification language is defined by the following BNF grammar:

\begin{small}
\begin{verbatim}
<spec>       ::= <safety_spec> | <liveness_spec> | <spec> "and" <spec>
<safety_spec> ::= "always" <predicate>
               |  "always" <predicate> "within" <time_bound>
<liveness_spec> ::= "eventually" <predicate>
               |  <predicate> "responds" <predicate> "within" <time_bound>
<predicate>  ::= <linear_pred> | <distance_pred> | <region_pred>
               |  <predicate> "and" <predicate>
               |  <predicate> "or" <predicate>
               |  "not" <predicate>
<linear_pred>   ::= <expr> <cmp_op> <expr>
<distance_pred> ::= "dist(" <agent_ref> "," <agent_ref> ")" <cmp_op> <num>
<region_pred>   ::= <agent_ref> "in" <region>
<agent_ref>     ::= "agent[" <int> "]" "." <field>
<field>         ::= "x" | "y" | "vx" | "vy" | "theta" | <string>
<region>        ::= "box(" <num> "," <num> "," <num> "," <num> ")"
               |  "circle(" <num> "," <num> "," <num> ")"
<expr>       ::= <num> | <agent_ref> | <expr> <arith_op> <expr>
<cmp_op>     ::= "<" | "<=" | ">" | ">=" | "=="
<arith_op>   ::= "+" | "-" | "*"
<time_bound> ::= <int> "steps"
<num>        ::= <float> | <int>
\end{verbatim}
\end{small}

\subsection{Denotational Semantics}

We define the semantics of specifications over infinite traces
$\tau = s_0, s_1, s_2, \ldots$ where each $s_t \in \calS$.

\begin{definition}[Trace satisfaction]
The satisfaction relation $\tau, t \models \psi$ is defined inductively:
\begin{align}
  \tau, t &\models \text{``always } \psi\text{''} &&\iff \forall t' \geq t.\; \tau, t' \models \psi \\
  \tau, t &\models \text{``eventually } \psi\text{''} &&\iff \exists t' \geq t.\; \tau, t' \models \psi \\
  \tau, t &\models \text{``}\psi_1 \text{ responds } \psi_2 \text{ within } k\text{''} &&\iff
    \tau, t \models \psi_1 \implies \exists t' \in [t, t+k].\; \tau, t' \models \psi_2 \\
  \tau, t &\models \psi_1 \text{ and } \psi_2 &&\iff \tau, t \models \psi_1 \text{ and } \tau, t \models \psi_2 \\
  \tau, t &\models \psi_1 \text{ or } \psi_2 &&\iff \tau, t \models \psi_1 \text{ or } \tau, t \models \psi_2 \\
  \tau, t &\models \text{not } \psi &&\iff \tau, t \not\models \psi
\end{align}
For atomic predicates:
\begin{align}
  \tau, t &\models e_1 \leq e_2 &&\iff \llbracket e_1 \rrbracket_{s_t} \leq \llbracket e_2 \rrbracket_{s_t} \\
  \tau, t &\models \text{dist}(i, j) > r &&\iff \norm{s_t^{(i)} - s_t^{(j)}} > r \\
  \tau, t &\models i \text{ in } R &&\iff s_t^{(i)} \in \llbracket R \rrbracket
\end{align}
where $\llbracket e \rrbracket_{s}$ denotes the evaluation of expression $e$
in state $s$, and $s_t^{(i)}$ is agent $i$'s state component at time $t$.
\end{definition}

\subsection{Robustness Semantics}

For quantitative analysis, we extend the Boolean semantics to a
\emph{robustness} semantics $\rho(\tau, t, \psi) \in \R$ following
the signal temporal logic (STL) tradition:
\begin{align}
  \rho(\tau, t, e_1 \leq e_2) &= \llbracket e_2 \rrbracket_{s_t} - \llbracket e_1 \rrbracket_{s_t} \\
  \rho(\tau, t, \text{always } \psi) &= \inf_{t' \geq t} \rho(\tau, t', \psi) \\
  \rho(\tau, t, \text{eventually } \psi) &= \sup_{t' \geq t} \rho(\tau, t', \psi) \\
  \rho(\tau, t, \psi_1 \text{ and } \psi_2) &= \min(\rho(\tau, t, \psi_1), \rho(\tau, t, \psi_2)) \\
  \rho(\tau, t, \text{not } \psi) &= -\rho(\tau, t, \psi)
\end{align}

The robustness value is positive when the specification is satisfied and negative
when violated. The safety margin reported by \marace{} (\Cref{def:margin})
corresponds to the robustness at the worst-case schedule.

\subsection{Compilation to Linear Constraints}

Safety specifications involving linear predicates compile directly to
the constraint form $Hs \leq h$ used by the abstract interpreter.
Distance predicates are linearized via the inscribed polytope
approximation: $\norm{x} > r$ is approximated by
$\bigvee_{k=1}^{K} v_k^T x > r$ where $\{v_k\}$ are uniformly spaced
unit vectors.

% ============================================================
\section{Probabilistic Estimation}
\label{app:prob}
% ============================================================

\subsection{Deployment Schedule Distributions}

In deployed MARL systems, the execution schedule is not adversarial but
follows a \emph{deployment distribution} $\mu_D$ determined by the
system's scheduling infrastructure (e.g., round-robin, event-driven,
or real-time scheduling). We model $\mu_D$ as a distribution over
HB-consistent schedules.

\begin{definition}[Race probability]
Given a deployment distribution $\mu_D$ over schedules, the \emph{race
probability} for an interaction race $(e, e', s_0)$ is:
\[
  p_R = \Prob_{\sigma \sim \mu_D}[\varphi(\mathrm{reach}(\sigma, s_0)) = 0
  \text{ and } e \text{ races with } e' \text{ under } \sigma].
\]
\end{definition}

\subsection{Importance Sampling Estimator}

Direct Monte Carlo estimation of $p_R$ is inefficient when $p_R$ is small
(\Cref{thm:stat-intract}). We use importance sampling with a proposal
distribution $q$ that concentrates mass near the race manifold:

\begin{definition}[IS estimator]
The self-normalized importance sampling estimator is:
\[
  \hat{p}_R^{\text{IS}} = \frac{\sum_{j=1}^{N} w_j \cdot \mathbf{1}[\text{race under } \sigma_j]}
  {\sum_{j=1}^{N} w_j},
  \quad w_j = \frac{\mu_D(\sigma_j)}{q(\sigma_j)},
  \quad \sigma_j \sim q.
\]
\end{definition}

\begin{theorem}[IS consistency]
\label{thm:is-consistent}
If $q(\sigma) > 0$ whenever $\mu_D(\sigma) > 0$ and $\sigma$ is in the
race set, then $\hat{p}_R^{\text{IS}} \xrightarrow{a.s.} p_R$ as $N \to \infty$.
\end{theorem}

\begin{proof}
By the strong law of large numbers applied to the numerator and denominator
separately. The numerator converges to $\E_q[w \cdot \mathbf{1}[\text{race}]]
= \E_{\mu_D}[\mathbf{1}[\text{race}]] = p_R$, and the denominator converges
to $\E_q[w] = 1$. Since both converge almost surely and the denominator is
bounded away from zero (as $q$ has full support on the $\mu_D$-support),
the ratio converges to $p_R$.
\end{proof}

\subsection{Concentration Bounds}

We provide finite-sample confidence intervals using the self-normalized bound:

\begin{theorem}[Self-normalized concentration]
\label{thm:concentration}
Let $w_1, \ldots, w_N$ be importance weights satisfying
$\max_j w_j / \bar{w} \leq M$ where
$\bar{w} = N^{-1}\sum_j w_j$.
Then for any $\delta > 0$:
\[
  \Prob\!\left[\abs{\hat{p}_R^{\text{IS}} - p_R} > \epsilon\right]
  \leq 2 \exp\!\left(-\frac{N_{\text{eff}} \epsilon^2}{2M^2}\right)
\]
where $N_{\text{eff}} = (\sum_j w_j)^2 / \sum_j w_j^2$ is the effective sample size.
\end{theorem}

\begin{proof}
The self-normalized estimator satisfies a Bernstein-type inequality.
Following Owen~\cite{owen2013monte}, the variance of $\hat{p}_R^{\text{IS}}$
is bounded by:
\[
  \Var[\hat{p}_R^{\text{IS}}] \leq \frac{p_R(1-p_R)}{N_{\text{eff}}}.
\]
The weight truncation $\max w_j / \bar{w} \leq M$ ensures sub-exponential
tails, and Bernstein's inequality yields the exponential concentration.
Specifically, define $Y_j = w_j(\mathbf{1}[\text{race under }\sigma_j] - p_R)/\bar{w}$.
Then $\E[Y_j] = 0$, $\abs{Y_j} \leq M$, and:
\[
  \Prob\!\left[\abs[\Big]{\sum_{j=1}^N Y_j} > N\epsilon\right]
  \leq 2\exp\!\left(-\frac{N^2\epsilon^2}{2\sum_j \E[Y_j^2] + 2NM\epsilon/3}\right).
\]
Since $\sum_j \E[Y_j^2] \leq NM^2/N_{\text{eff}} \cdot N$, simplification
gives the claimed bound.
\end{proof}

In our experiments, the effective sample size $N_{\text{eff}}$ equals 50
(the full sample count), indicating that the proposal distribution $q$
closely matches the target $\mu_D$ restricted to the race set, yielding
tight confidence intervals $[1.0, 1.0]$ for the planted races.

\subsection{Cross-Entropy Convergence}
\label{app:ce-convergence}

The proposal distribution $q$ is optimized via the cross-entropy (CE)
method~\cite{rubinstein2004cross}. Starting from a uniform proposal $q_0$
over HB-consistent schedules, the CE method iteratively:
\begin{enumerate}[nosep]
\item Draws $N$ samples from $q_t$;
\item Selects the top-$\rho$ fraction (``elite samples'') by race severity;
\item Fits $q_{t+1}$ to the elite samples by minimizing KL divergence.
\end{enumerate}

\begin{proposition}[CE convergence]
\label{prop:ce-conv}
Under regularity conditions (the parametric family of proposals includes
the optimal proposal $q^*$, and the elite threshold is chosen adaptively),
the CE iterates satisfy $D_{\text{KL}}(q^* \| q_t) \to 0$ at rate
$O(1/\sqrt{t})$.
\end{proposition}

The proof follows the standard CE convergence theory~\cite{rubinstein2004cross}
adapted to the constraint that proposals must assign zero mass to
HB-inconsistent schedules.


% ============================================================
\section{Lipschitz Analysis}
\label{app:lipschitz}
% ============================================================

The Lipschitz constant of agent policies plays a central role in both
the statistical intractability result (\Cref{thm:stat-intract}) and the
precision of the abstract domain. We describe the methods used to estimate
and bound Lipschitz constants.

\subsection{Spectral Norm Bounds}

For a ReLU network $\pi = W_L \circ \sigma \circ W_{L-1} \circ \cdots \circ \sigma \circ W_1$
where $\sigma$ is the ReLU activation, the global Lipschitz constant satisfies:
\[
  \Lip(\pi) \leq \prod_{\ell=1}^{L} \norm{W_\ell}_2.
\]

This spectral norm product bound is easy to compute but often loose.
We compute $\norm{W_\ell}_2$ via power iteration (10 iterations suffice
for convergence to relative error $< 10^{-6}$ in all our experiments).

\subsection{LipSDP Bounds}

For tighter bounds, we use the LipSDP framework~\cite{fazlyab2019efficient},
which formulates the Lipschitz constant computation as a semidefinite program:

\begin{align}
  \min_{\gamma, T} \quad & \gamma^2 \\
  \text{s.t.} \quad &
  \begin{bmatrix} W_1^T T W_1 - \gamma^2 I & W_1^T T \\ T W_1 & -T + W_2^T W_2 \end{bmatrix}
  \preceq 0, \\
  & T = \diag(t_1, \ldots, t_W), \quad t_i \geq 0.
\end{align}

This yields the tightest known polynomial-time computable bound on $\Lip(\pi)$
for ReLU networks.

\subsection{Lipschitz Calibration}

In \marace{}, the Lipschitz constant is used to calibrate the abstraction
error. Given $\Lip(\pi_i)$ for each agent, the maximum state perturbation
from a schedule change of magnitude $\Delta t$ is bounded by:
\[
  \norm{\Delta s_i} \leq \Lip(\pi_i) \cdot \Lip(T) \cdot \Delta t
\]
where $\Lip(T)$ is the Lipschitz constant of the transition function.

\begin{proposition}[Calibration convergence]
\label{prop:calibration}
As the number of calibration traces $K \to \infty$, the empirical
Lipschitz estimate $\hat{L}_K$ satisfies:
\[
  \hat{L}_K \leq \Lip(\pi) \leq \hat{L}_K + O\!\left(\sqrt{\frac{\log K}{K}}\right)
\]
with probability at least $1 - 1/K$.
\end{proposition}

\begin{proof}
The empirical Lipschitz estimate is:
\[
  \hat{L}_K = \max_{1 \leq j < k \leq K}
  \frac{\norm{\pi(s_j) - \pi(s_k)}}{\norm{s_j - s_k}}.
\]
This is a lower bound on $\Lip(\pi)$ by definition. The upper bound follows
from a covering number argument: with $K$ samples from a compact domain,
the uncovered volume has measure $O(K^{-1/d})$, and the maximum Lipschitz
ratio over the uncovered region is bounded by the global Lipschitz constant
plus a term that decreases as $O(\sqrt{\log K / K})$ by a uniform convergence
argument over Lipschitz functions.
\end{proof}

\subsection{Abstract Error Propagation}

The abstraction error $\epsilon_\ell$ at layer $\ell$ of the zonotope
analysis satisfies the recurrence:
\[
  \epsilon_{\ell+1} \leq \Lip(f_\ell) \cdot \epsilon_\ell + \delta_\ell
\]
where $\delta_\ell$ is the local approximation error introduced by the
ReLU relaxation at layer $\ell$. For the parallelotope relaxation,
$\delta_\ell \leq \sum_{k: l_k < 0 < u_k} \abs{u_k l_k} / (2(u_k - l_k))$.

The total abstraction error after $L$ layers is bounded by:
\[
  \epsilon_L \leq \epsilon_0 \prod_{\ell=0}^{L-1} \Lip(f_\ell)
  + \sum_{\ell=0}^{L-1} \delta_\ell \prod_{j=\ell+1}^{L-1} \Lip(f_j).
\]

In our experiments, the measured abstraction errors range from 0.036
(2-agent highway) to 0.121 (8-agent warehouse), consistent with the
theoretical bounds.


% ============================================================
\section{Compositional Decomposition}
\label{app:composition}
% ============================================================

\subsection{Assume-Guarantee Framework}

Compositional verification decomposes the $n$-agent system into $G$
interaction groups $\calG_1, \ldots, \calG_G$ where $\bigcup_g \calG_g = \calA$
and agents in different groups have limited interaction.

\begin{definition}[Interaction group contract]
An \emph{interaction group contract} for group $\calG_g$ is a pair
$(\mathsf{Assume}_g, \mathsf{Guarantee}_g)$ where:
\begin{itemize}[nosep]
\item $\mathsf{Assume}_g \subseteq \calS$ specifies the assumed behavior of
agents outside $\calG_g$ (their states remain within $\mathsf{Assume}_g$);
\item $\mathsf{Guarantee}_g \subseteq \calS$ specifies the guaranteed behavior
of agents within $\calG_g$ (their states remain within $\mathsf{Guarantee}_g$).
\end{itemize}
\end{definition}

\begin{theorem}[Assume-guarantee soundness]
\label{thm:ag-sound}
If for each group $g$:
\begin{enumerate}[nosep,label=(\roman*)]
\item Under assumption $\mathsf{Assume}_g$, the group $\calG_g$ satisfies
$\mathsf{Guarantee}_g$; and
\item $\mathsf{Guarantee}_h \subseteq \mathsf{Assume}_g$ for all groups $h \neq g$;
\end{enumerate}
then the full system satisfies $\bigcap_g \mathsf{Guarantee}_g$.
\end{theorem}

\begin{proof}
By induction on the number of execution steps. At step $t=0$, all agents are
in their initial states, which satisfy all assumptions by hypothesis.

Assume at step $t$, all agents satisfy their group guarantees. For each
group $g$, the agents outside $\calG_g$ are in $\bigcup_{h \neq g} \mathsf{Guarantee}_h
\subseteq \mathsf{Assume}_g$ by condition (ii). Therefore, by condition (i),
the agents in $\calG_g$ at step $t+1$ satisfy $\mathsf{Guarantee}_g$.

By induction, all agents satisfy their guarantees at all steps.
\end{proof}

\subsection{Group Size Theory}

The effectiveness of decomposition depends on the maximum interaction group
size, which determines the dimensionality of the per-group verification problem.

\begin{definition}[Interaction graph]
The \emph{interaction graph} $I = (\calA, E_I)$ has an edge $(i,j) \in E_I$
if agents $i$ and $j$ have concurrent events that could affect a shared
safety predicate. The interaction groups are the connected components of $I$.
\end{definition}

\begin{theorem}[Group size bound]
\label{thm:group-size}
If the interaction graph has maximum degree $\Delta$ and the safety
predicates have locality radius $r$ (each predicate involves agents within
distance $r$ in the interaction graph), then the maximum interaction group
size is bounded by $O(\Delta^r)$.
\end{theorem}

\begin{proof}
The interaction group containing agent $i$ is the connected component
reachable from $i$ in the interaction graph restricted to agents
participating in shared safety predicates. From any vertex in a graph
of maximum degree $\Delta$, the number of vertices reachable within
distance $r$ is at most $1 + \Delta + \Delta^2 + \cdots + \Delta^r
= O(\Delta^r)$.
\end{proof}

In spatially distributed multi-agent systems, $\Delta$ and $r$ are typically
small constants (each agent interacts with a bounded number of neighbors),
yielding constant-size groups and enabling the linear scaling observed in
our experiments.

\subsection{Percolation Bounds}

For randomly distributed agents, we can obtain probabilistic bounds on
group sizes using percolation theory.

\begin{proposition}[Sub-critical percolation]
\label{prop:percolation}
If $n$ agents are distributed uniformly at random in $[0,1]^2$ with
interaction radius $r_n = c \sqrt{\log n / n}$ for $c < 1/\sqrt{\pi}$,
then with probability $1 - O(1/n)$, the maximum interaction group size
is $O(\log n)$.
\end{proposition}

\begin{proof}
This follows from the theory of random geometric graphs. The expected
degree is $\lambda = n \pi r_n^2 = \pi c^2 \log n$. For $\pi c^2 < 1$,
the random geometric graph is in the sub-critical regime, where the
largest connected component has size $O(\log n)$ with high probability.
The result follows from the standard percolation threshold for
two-dimensional Poisson Boolean models, adapted to the finite setting
via a union bound over $n$ agents.
\end{proof}

This theoretical bound is consistent with our experimental observation
that all agents form a single interaction group in our benchmarks:
the highway intersection and warehouse corridor scenarios have
interaction radii above the critical threshold, as agents are densely
packed in shared spaces.


% ============================================================
\section{Proof Certificates}
\label{app:certificates}
% ============================================================

\marace{} can produce machine-checkable proof certificates that
independently verify the absence of interaction races.

\subsection{Certificate Format}

A proof certificate $\calP$ for race absence consists of:

\begin{definition}[Proof certificate]
\label{def:certificate}
$\calP = (\calZ_0, \{f^\sharp_\ell\}_{\ell=1}^{L_{\text{total}}}, \calC_{\text{HB}}, m^*)$
where:
\begin{itemize}[nosep]
\item $\calZ_0 = \zo(c_0, G_0)$ is the initial zonotope covering the
analyzed state region;
\item $\{f^\sharp_\ell\}$ is the sequence of abstract transfer functions
(linear maps and ReLU relaxation coefficients);
\item $\calC_{\text{HB}}$ is the set of HB constraints on generator
coefficients;
\item $m^* > 0$ is the certified minimum safety margin.
\end{itemize}
\end{definition}

\subsection{Verification Algorithm}

The certificate is verified by an independent checker that replays the
abstract computation and confirms the safety margin:

\begin{algorithm}[H]
\caption{Certificate Verification}
\label{alg:verify}
\begin{algorithmic}[1]
\Require Certificate $\calP = (\calZ_0, \{f^\sharp_\ell\}, \calC_{\text{HB}}, m^*)$,
safety predicate $(H, h)$
\Ensure Accept or Reject
\State $Z \leftarrow \calZ_0$; $\calC \leftarrow \calC_{\text{HB}}$
\For{$\ell = 1, \ldots, L_{\text{total}}$}
  \State $Z \leftarrow f^\sharp_\ell(Z)$ \Comment{Apply transfer function}
  \State Verify $f^\sharp_\ell$ is a sound abstraction of $f_\ell$
\EndFor
\State Compute $m_{\text{verified}} \leftarrow \min_k \min_{\xi \in [-1,1]^p, \calC}
  (h_k - H_k(c + G\xi)) / \norm{H_k}$
\If{$m_{\text{verified}} \geq m^*$ and $m^* > 0$}
  \State \Return \textbf{Accept}
\Else
  \State \Return \textbf{Reject}
\EndIf
\end{algorithmic}
\end{algorithm}

\begin{proposition}[Certificate soundness]
If \Cref{alg:verify} accepts certificate $\calP$, then the multi-agent
system has no interaction race for the specified safety predicate and
initial state region.
\end{proposition}

\begin{proof}
The verifier independently confirms that:
(1)~each transfer function is a sound over-approximation of the
corresponding concrete operation;
(2)~the HB constraints are consistent with the declared
happens-before graph;
(3)~the minimum safety margin over all feasible states in the
final zonotope is positive.
If all three hold, no concrete reachable state violates the
safety predicate.
\end{proof}

\subsection{Certificate Size}

The certificate size is polynomial in the system parameters:
$O(n \cdot L \cdot W^2 + n^2 \cdot T)$ where the first term accounts
for the weight matrices and relaxation coefficients, and the second term
accounts for the HB constraints. For our experimental configurations,
certificates range from approximately 2KB (2-agent highway) to 15KB
(8-agent warehouse).


% ============================================================
\section{Additional Experimental Details}
\label{app:experiments}
% ============================================================

\subsection{Policy Architectures}

\paragraph{Highway intersection.}
Each vehicle agent uses a 2-layer ReLU network: input dimension 4
(position $x, y$, velocity $v_x, v_y$), hidden layer of 64 units,
output dimension 2 (acceleration, steering). Total parameters: 450.

\paragraph{Warehouse corridor.}
Each robot agent uses a 2-layer ReLU network: input dimension 6
(position $x, y$, velocity $v_x, v_y$, orientation $\theta$, load status),
hidden layer of 128 units, output dimension 3 (forward velocity, angular
velocity, gripper). Total parameters: 1,155.

\subsection{Training Details}

Policies are trained using Proximal Policy Optimization (PPO)~\cite{schulman2017proximal}
with independent learning (each agent optimizes its own reward). Training
hyperparameters:
\begin{itemize}[nosep]
\item Learning rate: $3 \times 10^{-4}$ with linear decay;
\item Discount factor: $\gamma = 0.99$;
\item GAE parameter: $\lambda = 0.95$;
\item Clip ratio: $\epsilon = 0.2$;
\item Training episodes: 10,000 per agent;
\item Batch size: 64 episodes.
\end{itemize}

The reward function for each agent combines a task reward (reaching the
goal position) with a safety penalty ($-100$ for collision) and a
coordination bonus ($+1$ for maintaining safe separation from other agents).

\subsection{Race Planting Procedure}

To create ground-truth interaction races for evaluation, we use the
following procedure:

\begin{enumerate}[nosep]
\item \textbf{Identify critical regions:} Determine spatial regions where
multiple agents' paths intersect (intersection center, corridor junctions).
\item \textbf{Construct race-inducing initial conditions:} Place agents at
positions and velocities such that simultaneous action produces collision
but sequential (HB-ordered) action allows safe passage.
\item \textbf{Verify ground truth:} Exhaustively enumerate all HB-consistent
schedules for the initial condition (feasible for $n \leq 4$) and confirm
that exactly one pair of events constitutes a race.
\item \textbf{Embed in trace set:} Include the race-inducing initial condition
among the $K = 5$ trace episodes, with the remaining 4 episodes drawn from
the normal training distribution.
\end{enumerate}

For the warehouse corridor benchmark with $n \geq 6$ agents, exhaustive
enumeration is infeasible, so ground truth is established by combining
manual analysis with targeted simulation (1000 random schedules for each
candidate race pair).

\subsection{Time Breakdown}

\Cref{tab:time-breakdown} provides a detailed breakdown of analysis time
for each pipeline stage.

\begin{table}[h]
\centering
\caption{Time breakdown by pipeline stage (seconds). Percentages show
fraction of total time.}
\label{tab:time-breakdown}
\begin{tabular}{@{}llccccccc@{}}
\toprule
& & \multicolumn{6}{c}{\textbf{Pipeline Stage}} & \\
\cmidrule(lr){3-8}
\textbf{Scenario} & $n$ & Trace & HB & Decomp & Abstract & MCTS & IS & \textbf{Total} \\
\midrule
Highway & 2 & 0.012 & 0.001 & 0.003 & 0.001 & 0.020 & 0.002 & 0.039 \\
Highway & 3 & 0.027 & 0.001 & 0.004 & 0.001 & 0.037 & 0.002 & 0.072 \\
Highway & 4 & 0.048 & 0.001 & 0.006 & 0.001 & 0.063 & 0.002 & 0.121 \\
\midrule
Warehouse & 4 & 0.053 & 0.008 & 0.004 & 0.001 & 0.066 & 0.002 & 0.134 \\
Warehouse & 6 & 0.134 & 0.002 & 0.006 & 0.001 & 0.136 & 0.004 & 0.283 \\
Warehouse & 8 & 0.277 & 0.002 & 0.010 & 0.002 & 0.224 & 0.004 & 0.519 \\
\bottomrule
\end{tabular}
\end{table}

The dominant costs are trace collection and MCTS search, which together
account for 85--96\% of total time across all configurations. The abstract
interpretation and decomposition stages are negligible ($<2$\% of total time),
validating the efficiency of the constrained zonotope domain.

\subsection{Zonotope Statistics}

\begin{table}[h]
\centering
\caption{Zonotope propagation statistics for the detection experiments.}
\label{tab:zonotope-stats}
\begin{tabular}{@{}llcccc@{}}
\toprule
\textbf{Scenario} & $n$ & \textbf{HB Events} & \textbf{HB Edges} & \textbf{Coverage (\%)} & \textbf{Abs.\ Error} \\
\midrule
Highway & 2 & 300 & 285 & 3.62 & 0.036 \\
Highway & 3 & 400 & 380 & 3.98 & 0.040 \\
Highway & 4 & 500 & 475 & 4.71 & 0.047 \\
\midrule
Warehouse & 4 & 500 & 475 & 2.09 & 0.047 \\
Warehouse & 6 & 700 & 665 & 3.21 & 0.071 \\
Warehouse & 8 & 900 & 855 & 5.45 & 0.121 \\
\bottomrule
\end{tabular}
\end{table}

The HB graphs are dense (95\% of event pairs are causally ordered),
reflecting the tight physical coupling in both benchmarks. The coverage
percentage (fraction of the state space covered by the zonotope) ranges
from 2.09\% to 5.45\%, indicating that the abstract domain effectively
focuses on the relevant portion of the state space. The abstraction error
grows sub-linearly with agent count, consistent with the theoretical
bound from \Cref{app:lipschitz}.


% ============================================================
\section{Zonotope Reduction}
\label{app:reduction}
% ============================================================

As the analysis propagates through ReLU layers, the number of zonotope
generators grows. To maintain computational tractability, we periodically
reduce the generator count.

\subsection{Reduction Strategy}

Given a zonotope $Z = \zo(c, G)$ with $p > p_{\max}$ generators, we
reduce to $p_{\max}$ generators using the following procedure:

\begin{enumerate}[nosep]
\item \textbf{Sort generators by norm:} order $g_1, \ldots, g_p$ such that
$\norm{g_1} \geq \norm{g_2} \geq \cdots \geq \norm{g_p}$.
\item \textbf{Keep top generators:} retain $g_1, \ldots, g_{p_{\max}-d}$.
\item \textbf{Merge remaining generators:} replace $g_{p_{\max}-d+1}, \ldots, g_p$
with $d$ axis-aligned generators forming the interval hull:
\[
  g'_k = \left(\sum_{j=p_{\max}-d+1}^{p} \abs{g_{jk}}\right) e_k,
  \quad k = 1, \ldots, d.
\]
\end{enumerate}

\subsection{Bounded Error Analysis}

\begin{theorem}[Reduction error bound]
\label{thm:reduction-error}
Let $Z = \zo(c, G)$ be a zonotope and $Z' = \zo(c, G')$ the reduced zonotope
obtained by the above procedure. Then:
\[
  Z \subseteq Z' \subseteq Z \oplus B_\infty(\epsilon_{\text{red}})
\]
where $\epsilon_{\text{red}} = \max_{k} \sum_{j=p_{\max}-d+1}^{p} \abs{g_{jk}}$
and $B_\infty(\epsilon)$ is the $\ell_\infty$-ball of radius $\epsilon$.
\end{theorem}

\begin{proof}
\emph{Containment $Z \subseteq Z'$.}
For any $s = c + \sum_{j} \xi_j g_j \in Z$ with $\xi \in [-1,1]^p$,
we construct $\xi' \in [-1,1]^{p_{\max}}$ such that $s = c + G'\xi'$.
For the retained generators, $\xi'_j = \xi_j$ for $j \leq p_{\max} - d$.
For the merged generators, $\xi'_{p_{\max}-d+k} = \sum_{j > p_{\max}-d} \xi_j g_{jk} / g'_k$
for each axis $k$. By the triangle inequality:
\[
  \abs{\xi'_{p_{\max}-d+k}} = \frac{\abs[\big]{\sum_j \xi_j g_{jk}}}{\sum_j \abs{g_{jk}}}
  \leq \frac{\sum_j \abs{g_{jk}}}{\sum_j \abs{g_{jk}}} = 1.
\]
Hence $s \in Z'$.

\emph{Over-approximation bound.}
For any $s' = c + G'\xi' \in Z'$, we can write
$s' = c + \sum_{j=1}^{p_{\max}-d} \xi'_j g_j + \sum_{k=1}^{d} \xi'_{p_{\max}-d+k} g'_k e_k$.
The last $d$ terms contribute at most $\abs{g'_k}$ in each dimension $k$,
while the corresponding terms in $Z$ contribute at most
$\sum_{j > p_{\max}-d} \abs{g_{jk}} = g'_k$. The over-approximation error
is thus exactly $\epsilon_{\text{red}} = \max_k g'_k$ per dimension.
\end{proof}

\subsection{Adaptive Reduction}

In practice, we use an adaptive reduction strategy:
\begin{itemize}[nosep]
\item \textbf{Threshold:} Reduction is triggered when $p > 2 p_{\max}$
to amortize the cost.
\item \textbf{PCA-based merging:} Instead of axis-aligned merging, we
compute the principal components of the small generators and merge
along the top-$d$ principal directions. This reduces the approximation
error when the small generators are correlated.
\item \textbf{Error tracking:} The cumulative reduction error is tracked
and reported as part of the abstract error (\Cref{tab:zonotope-stats}).
\end{itemize}

\begin{proposition}[PCA merging improvement]
\label{prop:pca-merge}
PCA-based merging reduces the over-approximation volume by a factor of at
most $\prod_{k=1}^{d} \sigma_k / \sigma_k'$ compared to axis-aligned merging,
where $\sigma_k$ and $\sigma_k'$ are the singular values of the small generator
matrix in the original and rotated coordinate systems, respectively.
\end{proposition}

\begin{proof}
The volume of the over-approximation box in the axis-aligned case is
$\prod_k 2\sum_j \abs{g_{jk}}$. In the PCA-rotated coordinate system, the
box volume is $\prod_k 2\sigma_k$ where $\sigma_k$ are the singular values
of $[g_{p_{\max}-d+1} \cdots g_p]$. Since the singular values satisfy
$\sigma_k \leq \sum_j \abs{g_{jk}}$ with equality only when all small
generators are axis-aligned, PCA merging always yields equal or smaller
over-approximation volume.
\end{proof}

For the configurations in our experiments, PCA merging was not necessary
(the generator count stayed below $p_{\max} = 100$ throughout), but the
mechanism is available for larger networks with more unstable neurons.

\end{document}
